\chapter{Adult-Onset Still's Disease}
\index{Adult-Onset Still's Disease}
\label{sec:Adult-Onset Still's Disease (AOSD)}

\section{Overview}
Adult-onset Still's disease is a rare systemic inflammatory disorder of unknown cause, considered the adult counterpart of systemic juvenile idiopathic arthritis (sJIA), characterized by a classic triad of high-spiking fevers, evanescent salmon-pink rash, and joint pain/arthritis, with an incidence of 0.16--0.4 per 100,000, a bimodal age distribution (15--25 and 36--46 years), and a slight female predominance.
\section{Causes}
This condition happens because of dysregulated innate immunity with activation of the inflammasome, leading to overproduction of IL-1, IL-6, IL-18, and other pro-inflammatory cytokines, with macrophage activation being a key feature. The underlying mechanisms involve complex interactions between genetic predisposition and environmental factors that disrupt normal physiological processes. The underlying mechanisms involve complex interactions between genetic predisposition and environmental factors. Disruption of normal physiological processes leads to the characteristic features of the condition. Multiple contributing factors may act synergistically to trigger or worsen the condition. Understanding the specific cause helps guide targeted treatment approaches.
\section{Risk Factors}


\begin{itemize}[noitemsep]
  \item Genetic predisposition and family history
  \item Advancing age
  \item Lifestyle factors (diet, exercise, smoking, alcohol)
  \item Environmental exposures
  \item Underlying medical conditions
  \item Medication use
\end{itemize}
\section{Signs and Symptoms}

\subsection{Common symptoms}
\begin{itemize}[noitemsep]
  \item You may experience quotidian or double-quotidian fevers, an evanescent, salmon-pink, non-pruritic maculopapular rash (Koebner phenomenon), joint pain and arthritis (initially pauciarticular, may progress to severe destructive polyarthritis, particularly involving wrists and carpal bones---carpal ankylosis is characteristic), sore throat (intense, exudative, a highly distinctive feature), lymphadenopathy and splenomegaly, serositis, and muscle pain. Symptoms vary depending on the severity and extent of the condition Pain or discomfort in the affected area Functional impairment affecting daily activities Changes in appearance or function of affected tissues Systemic symptoms may include fatigue, fever, or weight changes Symptoms may develop gradually or appear suddenly
  \item Pain or discomfort in the affected area
  \end{itemize}

\subsection{Emergency warning signs}
\begin{itemize}[noitemsep]
  \item Severe or worsening symptoms of adult-onset still's disease require immediate medical evaluation
\end{itemize}
\section{Progression and Stages}
The condition may progress through different stages of severity over time. Early stages often involve mild or subtle symptoms that may be overlooked. Moderate stages involve more noticeable symptoms affecting daily function. Advanced stages may involve significant impairment and complications.
\section{Complications}
\begin{itemize}[noitemsep]
  \item Disease progression leading to worsening symptoms
  \item Secondary infections or injuries
  \item Side effects of treatment
  \item Reduced quality of life
  \item Emotional and psychological impact
\end{itemize}
\section{Diagnosis}
Doctors diagnose this using Yamaguchi criteria, with serum ferritin characteristically markedly elevated (often $>$5--10 times normal, $>$1000 ng/mL), with decreased glycosylated ferritin ($<$20\%). Comprehensive medical history and physical examination Laboratory tests to assess organ function and rule out other conditions Imaging studies to visualize affected structures Specialized testing depending on the affected system Consultation with relevant specialists Monitoring of disease progression over time

\begin{itemize}[noitemsep]
  \item Comprehensive medical history and physical examination
  \item Laboratory tests to assess organ function
  \item Imaging studies to visualize affected structures
  \item Specialized testing depending on the affected system
  \item Consultation with relevant specialists
\end{itemize}
\section{Prevention}
\begin{itemize}[noitemsep]
  \item Maintain a healthy lifestyle with balanced diet and exercise
  \item Avoid known risk factors
  \item Attend regular medical check-ups for early detection
  \item Follow recommended screening guidelines
\end{itemize}
\section{Treatment Options}
Treatment typically involves stratified: NSAIDs for mild disease, corticosteroids for moderate disease, and IL-1 inhibitors (anakinra, canakinumab) or IL-6 inhibitors (tocilizumab) for refractory or severe disease. Medications to manage symptoms and address underlying mechanisms Lifestyle modifications and self-care measures Physical therapy and rehabilitation Surgical intervention when indicated Regular monitoring and follow-up care Patient education and support services

\begin{itemize}[noitemsep]
  \item Medications to manage symptoms and address underlying mechanisms
  \item Lifestyle modifications and self-care measures
  \item Physical therapy and rehabilitation
  \item Surgical intervention when indicated
  \item Regular monitoring and follow-up care
\end{itemize}
\section{Living with It}
\begin{itemize}[noitemsep]
  \item Follow your treatment plan as prescribed
  \item Maintain regular follow-up with your healthcare provider
  \item Make necessary lifestyle adjustments
  \item Seek support from family, friends, or support groups
  \item Stay informed about your condition
\end{itemize}
\section{Outlook}
\begin{itemize}[noitemsep]
  \item The outlook varies from person to person
  \item Some patients achieve remission, while others have a chronic relapsing course.
\end{itemize}
